\subsection{Emotion's Function in the Chain}

The chapter has settled on a coordinated notation for the actualization chain: nodes $A_0$ through $A_4$ for the major actualization stages, with $A_3$ as the emotion node (cf. \S 1.7); and transitions $M_n \rightarrow M_{n+1}$ for the moments between nodes, with $M_3 \rightarrow M_4$ designating the conversion that emotion \textit{is}. What the chain requires at $M_3 \rightarrow M_4$, under evaluatively complex conditions, is the conversion of \textit{orexis} (\gk{ὄρεξις} ``desire'') into a determinate form of desire directed at a specific object under a specific evaluative description. Emotion is precisely this conversion. Without it, the chain would have a generic moved mover---\textit{orexis} set in motion by the apparent good (\textit{De Anima} III.10, 433a31-b1)---but no mechanism to specify \textit{which} form of desire corresponds to \textit{which} disclosure of the apparent good. The desire to flee is not the desire to retaliate; the desire to relieve is not the desire to conceal. Each determinate conative orientation is individuated by the cognitive structure of the emotion that specifies it (cf. \S 1.4), and each draws upon a determinate pattern of somatic preparation---the cooling of fear, the heating of anger, the sympathetic pain of pity---that configures the body for the action the emotion calls for (cf. \S 1.5).

Emotion does not precede desire as a separate antecedent stage; emotion \textit{is} the form desire takes when \textit{doxa} has disclosed evaluatively complex content. Fear is not something that happens before the desire to flee; fear \textit{is} the desire to flee insofar as it is structured by the cognitive evaluation of an anticipated destructive evil and experienced in the somatic alteration of chilled blood. Anger is not something that happens before the desire to retaliate; anger \textit{is} the desire to retaliate insofar as it is structured by the perception of an unjustified slight and experienced in the somatic alteration of heated blood. Pity is not a separate state that subsequently produces the wish that undeserved suffering cease; pity is the painful recognition of undeserved suffering inseparable from that wish---and similarly for shame, kindness, indignation, and the rest of the \textit{Rhetoric}'s catalog. The \textit{Rhetoric}'s definitions confirm this composite structure (cf. \S 1.4): each emotion is a desire of a certain kind, attended by pain or pleasure, directed at a certain kind of object, under a certain evaluative description. The composite is not several parallel components stacked together but a single articulated determination in which the cognitive evaluation, the conative orientation, the hedonic tonality, and the physiological alteration are mutually constitutive aspects of one and the same \textit{pathos} (cf. \S 1.3). What \S 1.3 established as the composite-structure of the enmattered \textit{pathos}, the present section locates in the actualization chain: the composite is precisely what is required at $M_3 \rightarrow M_4$ in the evaluatively complex case, and emotion is precisely the composite operative as conversion.

The three-factor schema of \textit{De Anima} III.10 (433a9-b30) is preserved without revision. \textit{Orexis} remains the moved mover, set in motion by the apparent good and itself moving the bodily instrument; the realizable good remains the unmoved originator, drawing the desiring soul without being drawn in turn; the animal via its bodily instrument remains that which is moved, traversing the field of action under the direction of desire. What the schema's original formulation leaves unspecified is the \textit{internal} complexity the moved mover acquires when the apparent good is articulated under the evaluative categories of the \textit{Rhetoric}---when the slighter, the threat, the undeserved sufferer, the public falling-short have entered the soul through \textit{doxa} as evaluatively determinate disclosures of the apparent good. Emotion names this internal complexity. It is the actualization of the \textit{orektikon} (\gk{ὀρεκτικόν} ``faculty of desire'') into a specific form of desire directed at a specific object under a specific evaluative description; it is the form in which the soul's being-moved becomes the body's being-moved; it is the mechanism through which the world's disclosure as mattering in a certain way becomes the animal's motion in the corresponding direction. The schema persists; the moved mover gains an articulational concretion the bare formulation does not require but does not preclude.

The conversion-function emotion supplies is therefore not universal across the chain but specific to the evaluatively complex case. Where the apparent good is straightforwardly pleasant or painful and the desire is basic appetite or aversion---the drinking case of \textit{Movement of Animals} (701a32-33), in which appetite, perception or imagination, and immediate action follow without engagement of the higher-order \textit{pathē}---\textit{resonant orexis} alone suffices to direct the organism toward or away from the object (cf. \S 1.6). At $M_3 \rightarrow M_4$ in the appetitive case, the chain runs without internal complexity: \textit{orexis} is generic, and the body's being-moved is the immediate consequence of the perception's hedonic valence. The higher-order conversion is reserved for the case in which \textit{doxa} has done its articulational work. The present section has located emotion in the chain; the conclusion turns to the five-fold ontological structure each higher-order \textit{pathos} exhibits when so located.
